\documentclass[11pt]{article}

\usepackage{series}
\usepackage{amsmath,amssymb,amsthm}
\usepackage{geometry}
\usepackage{hyperref}
\usepackage{booktabs}

\geometry{margin=1in}

% --- Metadata ---
\title{Theta Closure in Modal Triplet Theory V:
Weak-Angle Round Trip and the Non-Circularity Criterion}
\author{Peter Nero}
\date{July 2026}

\begin{document}
\maketitle

\begin{abstract}
We reassess whether the weak mixing angle provides a redundant test of the
selected MTT gauge profile. Let
$r_{21}=I_2/I_1=g_1^2/g_2^2$, with $g_1=\sqrt{5/3}\,g_Y$. Then
$\sin^2\theta_W=3r_{21}/(5+3r_{21})$ at the same scale and in the same scheme.
For the selected SMDR profile at $Q=M_t$, this gives
$\sin^2\theta_W=0.2346644\pm0.0000433$. Because $r_{21}$ was itself obtained
from the measured $(g_Y,g_2)$ profile, this equality is an exact algebraic
round trip, not a held-out prediction. Supplying an absolute $g_2$ from
$(G_F,m_W)$ does not remove that circularity. A genuine redundancy test
requires an MTT source theorem selecting $r_{21}$ without electroweak gauge
data, followed by independent common-scheme transport. The obsolete one-loop
$5~\mathrm{TeV}$ calculation and its precision-prediction claim are withdrawn.
\end{abstract}

\section*{Revision note for this edition}
\addcontentsline{toc}{section}{Revision note for this edition}
\begin{description}
\item[Supersedes.] \emph{Theta Closure in Modal Triplet Theory V: Redundant
Determination from Gauge Couplings and the Weak Mixing Angle}, first edition.
\item[Reason.] The weak angle was presented as an independent precision test
even though it was algebraically reconstructed from the same measured gauge
profile, and the scale used the retired few-TeV crossing.
\item[Resolution.] Version~2 derives the same-scheme round-trip formula,
propagates its uncertainty at $Q=M_t$, and states an explicit non-circularity
criterion for any future test.
\item[Retained result.] The weak angle remains an exact consistency identity
for a supplied gauge-ratio profile.
\item[Remaining boundary.] A held-out test requires MTT to select $r_{21}$
without electroweak gauge data and then transport it independently.
\end{description}

\section{Selected inputs and scope}

Paper~I adopts the full-Standard-Model $\overline{\mathrm{MS}}$ profile
transported by SMDR~v1.3 to
\[
Q=M_t=172.5590883453979~\mathrm{GeV}.
\]
Its relevant entries are
\[
g_Y=0.3585945042\pm0.0000307251,
\qquad
g_2=0.6475986708\pm0.0000287665,
\]
or $g_1=\sqrt{5/3}\,g_Y$ in GUT normalization. The associated overlap ratio is
\[
r_{21}:=\frac{I_2}{I_1}=0.5110273\pm0.0001231.
\]
Paper~II realizes this ratio by calibrating its effective geometry. Paper~III
checks conditional compatibility of the quadratic norm. Neither paper
currently supplies a value-source theorem selecting $r_{21}$ independently of
the gauge profile.

\section{Same-scale weak-angle identity}

The overlap convention is
\[
\frac{1}{g_a^2}=\frac{I_a}{g_{10}^2}.
\]
Therefore
\begin{equation}
r_{21}=\frac{I_2}{I_1}=\frac{g_1^2}{g_2^2}.
\label{eq:r21}
\end{equation}
Using $g_Y^2=(3/5)g_1^2$ and the $\overline{\mathrm{MS}}$ definition
\[
s_W^2(Q):=\frac{g_Y^2(Q)}{g_Y^2(Q)+g_2^2(Q)},
\]
gives the exact identity
\begin{equation}
\boxed{s_W^2(Q)=\frac{3r_{21}(Q)}{5+3r_{21}(Q)}.}
\label{eq:weak_identity}
\end{equation}

\begin{theorem}[Gauge-profile round-trip theorem]
If $r_{21}(Q)$ is constructed from the same measured common-scheme pair
$(g_Y(Q),g_2(Q))$ through Equation~\eqref{eq:r21}, then evaluating
Equation~\eqref{eq:weak_identity} returns the weak angle encoded in that pair.
The equality is algebraic and cannot constitute an independent prediction.
\end{theorem}

\begin{proof}
Substitution of $r_{21}=(5/3)g_Y^2/g_2^2$ into
Equation~\eqref{eq:weak_identity} gives
$g_Y^2/(g_Y^2+g_2^2)$ identically.
\end{proof}

For the selected ratio,
\begin{equation}
s_W^2(M_t)=0.2346644\pm0.0000433,
\end{equation}
where the uncertainty is propagated from the selected $r_{21}$ row. Direct
substitution of the selected $g_Y$ and $g_2$ gives the same central value. This
is a useful convention and arithmetic check, but it has no held-out status.

\section{Why the $(G_F,m_W)$ construction remains circular}

An electroweak input such as $(G_F,m_W)$ can set an absolute $SU(2)$
normalization, subject to radiative matching. It cannot make the weak-angle
test independent when the ratio $r_{21}$ still comes from the measured
hypercharge-to-$SU(2)$ profile. Indeed, once $g_2$ and $r_{21}$ are supplied,
Equation~\eqref{eq:r21} defines $g_1$ and hence the weak angle. The information
being tested is already present in $r_{21}$.

The earlier tree-level value near $0.2312$ used the obsolete $5~\mathrm{TeV}$
profile and a one-loop return to $M_Z$. It is withdrawn as a prediction. Its
threshold scan also showed that the apparent precision was not stable under
the stated electroweak matching variation.

\section{Criterion for a genuine redundant determination}

A non-circular weak-angle test requires all of the following:
\begin{enumerate}
\item MTT geometry selects $r_{21}$ without using $g_Y$, $g_1$, $g_2$, or
$\sin^2\theta_W$ as value inputs;
\item an absolute gauge normalization and matching scale are selected without
using the held-out weak angle;
\item the boundary data are transported with a declared multi-loop scheme and
threshold prescription; and
\item the resulting $s_W^2$ is compared with a measurement not used anywhere
in selection, calibration, or branch choice.
\end{enumerate}
Only then is Equation~\eqref{eq:weak_identity} a prediction map rather than a
round-trip map. Current Papers~I--III satisfy the algebraic and realization
parts but not the first value-source condition.

\section{Conclusion}

The weak mixing angle is exactly consistent with the selected gauge profile,
as it must be. The precise same-scale result is
$s_W^2(M_t)=0.2346644\pm0.0000433$. This validates the hypercharge
normalization and overlap-ratio bookkeeping.

It does not add an independent Standard Model observable to MTT closure. The
former non-circularity claim fails because the overlap ratio already contains
the measured weak-angle information. The next theorem target is sharply
defined: select $r_{21}$ from MTT source geometry before consulting the
electroweak gauge profile, then execute a held-out common-scheme comparison.


% ===========================
\section*{References}
% ===========================

\begin{enumerate}

\item Particle Data Group (PDG),
\emph{Review of Particle Physics},
Prog. Theor. Exp. Phys. 2024, 083C01
(for $G_F$, $m_W$, $M_Z$, and electroweak input values).

\item M. E. Machacek and M. T. Vaughn,
``Two-loop renormalization group equations in a general quantum field theory,''
Nucl. Phys. B222 (1983) 83--103.

\item T. Kato,
\emph{Perturbation Theory for Linear Operators},
Springer (1995).

\item P. Nero,
\emph{Modal Triplet Theory: Foundation},
MTT corpus.

\item P. Nero,
\emph{Modal Triplet Theory: Quantum Amplitudes from Modal Geometry},
MTT corpus.

\item P. Nero,
\emph{Gauge Couplings, Internal Geometry, and $\Theta$--Closure in Modal Triplet Theory},
Paper~I.

\item P. Nero,
\emph{Direct Geometric Evaluation of Nonabelian Overlaps in Modal Triplet Theory},
Paper~II.

\item P. Nero,
\emph{Twistor--Action Matching of Gauge Overlaps in Modal Triplet Theory},
Paper~III.

\item P. Nero,
\emph{Extending $\Theta$--Closure to Gravity and Cosmology in Modal Triplet Theory},
Paper~IV.



\end{enumerate}
\begin{thebibliography}{99}

\bibitem{MTT-Admissibility}
P.~Nero,
\emph{Modal Triplet Theory: Admissibility, Encodings, and the Structure of Physical Description},
Zenodo preprint, January 2026.
\url{https://doi.org/10.5281/zenodo.18255621}

\bibitem{MTT-Foundation}
P.~Nero,
\emph{Modal Triplet Theory: Foundation},
Zenodo preprint, September 2025.
\url{https://doi.org/10.5281/zenodo.16949762}

\bibitem{MTT-FixedPoints}
P.~Nero,
\emph{Fixed Points I--VI: Complete Coherence Spine},
Zenodo preprints, August 2025.
\url{https://doi.org/10.5281/zenodo.16948748}

\bibitem{MTT-ProjectionAdmissibility}
P.~Nero,
\emph{The Projection--Admissibility Principle: Structural Constraints on Effective Physical Description},
Zenodo preprint, January 2026.
\url{https://doi.org/10.5281/zenodo.18255838}

\bibitem{MTT-Closure}
P.~Nero,
\emph{Closure and Inevitability in Modal Triplet Theory},
Zenodo preprint, January 2026.
\url{https://doi.org/10.5281/zenodo.18255510}

\bibitem{MTT-CoherenceCapacity}
P.~Nero,
\emph{Coherence Capacity as the Fundamental Resource of Effective Physics},
Zenodo preprint, January 2026.
\url{https://doi.org/10.5281/zenodo.18255905}

\bibitem{MTT-CoherenceDynamics}
P.~Nero,
\emph{Dynamics of Coherence Capacity: Transport, Concentration, and Exhaustion},
Zenodo preprint, January 2026.
\url{https://doi.org/10.5281/zenodo.18256048}

\bibitem{MTT-QM}
P.~Nero,
\emph{Modal Triplet Theory: From MTT to Quantum Mechanics},
Zenodo preprint, September 2025.
\url{https://doi.org/10.5281/zenodo.17074246}

\bibitem{MTT-QFT}
P.~Nero,
\emph{From MTT to Quantum Field Theory},
Zenodo preprint, 2025.
\url{https://doi.org/10.5281/zenodo.17068816}

\bibitem{MTT-GR}
P.~Nero,
\emph{Modal Triplet Theory: From MTT to General Relativity},
Zenodo preprint, October 2025.
\url{https://doi.org/10.5281/zenodo.16950597}

\bibitem{MTT-UVQG}
P.~Nero,
\emph{Modal Triplet Theory: From MTT to a UV-Finite, Unitary Quantum Gravity},
Zenodo preprint, 2025.
\url{https://doi.org/10.5281/zenodo.17077671}

\bibitem{MTT-Measurement}
P.~Nero,
\emph{Measurement as Disturbance and Stabilization in Modal Triplet Theory},
Zenodo preprint, 2025.
\url{https://doi.org/10.5281/zenodo.17177404}

\bibitem{MTT-Irreversibility}
P.~Nero,
\emph{Projection, Probability, and Irreversibility: Shadow Bridges Between Measurement, Black Holes, and Cosmology in Modal Triplet Theory},
Zenodo preprint, January 2026.
\url{https://doi.org/10.5281/zenodo.18256408}

\bibitem{MTT-Bell-Beables}
P.~Nero,
\emph{Modal Fixed Points, Bell’s Beables, and the Limits of Factorization},
Zenodo preprint, 2025.
\url{https://doi.org/10.5281/zenodo.17076300}

\bibitem{MTT-Bell-Temporal}
P.~Nero,
\emph{Temporal Bell Inequalities and Global Consistency in Modal Triplet Theory},
Zenodo preprint, August 2025.
\url{https://doi.org/10.5281/zenodo.18208884}

\bibitem{MTT-Stochastic}
P.~Nero,
\emph{From Modal Triplet Theory to Indivisible Stochastic Processes: A First-Principles, Fully Rigorous Derivation},
Zenodo preprint, January 2026.
\url{https://doi.org/10.5281/zenodo.18254862}

\end{thebibliography}
\end{document}
